\section{Contractions with different boundary data}
\label{g:ba}

Theorem~\ref{g:cf:contraction-theorem} keeps a generic section and its
adjoint rank. The boundary constructions of Section~\ref{g:bt} also accommodate
three other useful forms of control: arbitrary exceptional twists,
eventual discrepancy inequalities, and pole tests on the original
variety.  We verify these forms separately.  This identifies the exact
support condition that must survive each dimension decrease.

Throughout this section $X$ is smooth connected projective,
$H^0(X,\Omega_X^q)=0$ for $q>0$, and $L=-K_X$ has a smooth semipositive
metric.  A torus $T$ acts with its natural linearization on $L$.
Every smooth projective rational image $Y$ has no positive-degree
holomorphic forms, so $\chi(\mathcal O_Y)=1$ and
$\operatorname{Pic}(Y)\otimes\mathbb Q$ is finite dimensional.
We use geometrically integral generic fibres throughout.

Two of the pole formulations below express the same uniqueness
condition. Given the distinguished generic section, the bound for
every exceptional twist in \eqref{g:bt:all-exceptional-rank} is equivalent
to original-prime uniqueness in Lemma~\ref{g:cf:uniqueness}.
Indeed, finitely many invariant rational sections regular at the
original horizontal primes acquire only exceptional horizontal poles
on a fixed equivariant resolution. One common effective invariant
exceptional divisor clears those poles. Conversely an exceptional
twist changes no order at an original prime. The constructions below
remain different: they use, respectively, a weighted metric and Iitaka
rigidity, or a signed boundary and the null-prime field.

\subsection{Rigid Iitaka fibres and arbitrary exceptional twists}
\label{g:ba:iitaka}

The following elementary form of Iitaka rigidity is useful because the
base may need further resolution after its Iitaka map is chosen.

\begin{lemma}\label{g:ba:rigidity}
Let $Q\geq0$ be an integral divisor on a smooth projective variety $Y$,
and let $u:Y'\to Y$, $h:Y'\to Y_2$ resolve its Iitaka fibration, with
connected generic fibre $S$.  For every effective $u$-exceptional
integral divisor $A$ and every positive integer $m$,
\[
             h^0\bigl(S,m(u^*Q+A)|_S\bigr)=1.
\]
\end{lemma}

\begin{proof}
Choose the Iitaka map from a multiple of $Q$ attaining the maximal image
dimension and take its Stein factorization.  A multiple of $u^*Q$
dominates the pullback of a big divisor on $Y_2$; after taking a higher
multiple and replacing that divisor by an ample divisor plus an
effective divisor, it dominates a pulled-back ample divisor.
The generic fibre is geometrically integral.
Adding an effective exceptional divisor changes no global space of
sections, by normality, and hence no Iitaka dimension.
If a multiple of $(u^*Q+A)|_S$ had two independent sections, extend
them rationally over $Y'$ and clear their vertical poles with a
sufficiently high pulled-back ample twist.  The preceding domination
absorbs that twist into a multiple of $u^*Q+A$.
Their ratio remains nonconstant on $S$.  Combined with ratios defining
the base, it increases the transcendence degree of a global section-ratio
field beyond $\kappa(Q)$, a contradiction.
The canonical effective divisor supplies the nonzero section.
\end{proof}

\begin{proposition}[Preserving arbitrary exceptional twists]
\label{g:ba:exceptional-contraction}
Let $\pi:Z\to X$, $g:Z\to Y$ be a smooth equivariant resolution
of a dominant $T$-invariant rational map to a smooth projective $Y$,
with $\dim Y>0$ and geometrically integral generic fibre
$F$ over $K=\mathbb C(Y)$. Suppose that
$H^0(F,a\pi^*L|_F)^T\ne0$ for some $a>0$, and that
\[
 \dim_K H^0(F,(b\pi^*L+A)|_F)^T\leq1
\]
for every integer $b\geq0$ and every effective invariant
$\pi$-exceptional integral divisor $A$.
Then either a positive multiple of $L$ has a nonzero invariant section
on $X$, or a dominant rational map $Y\dashrightarrow Y_2$ to a smooth
projective base of smaller dimension preserves both the generic section
and this bound for every exceptional twist on a smooth equivariant
resolution of the composite map.
\end{proposition}

\begin{proof}
Use the divisors $N,R,B$ of Proposition~\ref{g:bt:weighted-transfer}.
If any $j(N+R)+lB$, $j>0$, $l\geq0$, is rationally effective, that
proposition gives the section on $X$.  Assume otherwise.
The weighted cotangent bound gives $hcB-c_1(M)$ pseudoeffective for a
line $M$ in a cotangent tensor power of order $h$.  Consequently
\[
                    M+j(N+R)+B
                    \quad\hbox{is not big}\qquad(j>0);
\]
adding the pseudoeffective class $hcB-M$ to a big class would contradict
the preceding noneffectivity.

The bounded metric of $N$, hard Lefschetz and the nonzero constant term
of $\chi(K_Y+mN)$ give, by the determinant construction
of Lemma~\ref{g:determinant}, a fixed cotangent line $M$ and
effective $D_j\sim M+jN$ for unbounded $j>0$.
Fix two indices $i<b$ and put $Q=D_b+bR+B$.  This divisor is effective
and not big.  Resolve its Iitaka map by $u:Y'\to Y$, $h:Y'\to Y_2$.
Then $\dim Y_2<\dim Y$.

Restrict pullbacks to the generic fibre $S$ of $h$.
The interpolation equivalence \eqref{g:cf:interpolation}, applied at
the twist exponents $i<b<j$, gives two effective representatives of
the class $(j-i)D_{b,S}$.  Add the same effective complement
$(j-i)(Q_S-D_{b,S})$ to both.  Lemma~\ref{g:ba:rigidity} makes the
resulting representatives of $(j-i)Q_S$ equal, and cancellation gives
\[
           (j-i)D_{b,S}=(b-i)D_{j,S}+(j-b)D_{i,S}
                                                  \qquad(j>b).
\]
In particular all the later $D_{j,S}$ are supported in $D_{b,S}$.
Since $j$ is unbounded and $D_{j,S}\geq0$, the divisor
\[
                   V_S=\frac{D_{b,S}-D_{i,S}}{b-i}
\]
is effective and rationally equivalent to $N_S$.
Its support lies in $D_{b,S}$.  For every effective $u$-exceptional $A$
and integers $n,l\geq0$ this also gives
\begin{equation}\label{g:ba:exceptional-fibre-bound}
 h^0\bigl(S,n(N+R)_S+lB_S+A|_S\bigr)\leq1.
\end{equation}
To see this, represent $N_S$ by $V_S$.  The indicated effective
representative is supported in $(u^*Q+A)|_S$ and therefore dominated
by a multiple of it.  Two independent sections would remain independent
after taking powers and multiplying by a fixed section, contrary to
Lemma~\ref{g:ba:rigidity}.

On a common resolution over $Y'$, lift a denominator-clearing multiple
of $V_S$ through $f^*N\to ap^*L$.  It gives the new invariant generic
section. For the rank assertion, let $F_{\mathrm{exc}}$ be any effective
invariant divisor exceptional over $X$ on the common resolution.
A section with this twist over the new generic fibre is, by the old
rank condition, rationally
$\lambda s^n s_{F_{\mathrm{exc}}}$ with $\lambda\in\mathbb C(Y)$, first for exponents
divisible by $a$.  At a prime of $S$ not exceptional over $Y$ and not
in $B_S$, choose a nonexceptional component realizing the restricted
minimum $n'_Q=(N+R)_Q$.  The exceptional twist has coefficient zero
there.  Regularity gives
\[
                       \operatorname{ord}_Q\lambda
                                  \geq-n(N+R)_Q.
\]
All remaining poles occur on $B_S$ or on exceptional divisors of $u$.
For any proposed pair of sections, choose an effective $u$-exceptional
divisor $A_Y$ on $Y'$ and $l\geq0$ so that
$lB_S+A_Y|_S$ absorbs them. Equation~\eqref{g:ba:exceptional-fibre-bound}
then forces dependence.  Other exponents follow by taking powers to
clear divisibility; independent sections cannot become dependent under
a power on a geometrically integral fibre.
\end{proof}

\begin{proposition}[The unweighted rigid-fibre variant]
\label{g:ba:toroidal-contraction}
Let $\pi:Z\to X$, $g:Z\to Y$ be a smooth equivariant resolution
of a dominant $T$-invariant rational map to a smooth projective $Y$,
with $\dim Y>0$ and geometrically integral generic fibre
$F$ over $K=\mathbb C(Y)$. Suppose there is a nonzero invariant section
of $a\pi^*L|_F$, $a>0$, and
\[
 \dim_K H^0(F,(K_{Z/X}+ma\pi^*L)|_F)^T=1
                    \qquad(m\geq0).
\]
Then either a positive multiple of $L$ has a nonzero invariant section
on $X$, or a dominant rational map $Y\dashrightarrow Y_2$ to a smooth
projective base of smaller dimension preserves the generic section and
these Jacobian-twisted ranks, with a possibly different positive
integer $a$.
\end{proposition}

\begin{proof}
Take the prepared toroidal model of
Proposition~\ref{g:bt:toroidal-transfer}.
Let $B$ now be the reduced sum of primes $Q$ for which no component
attaining the minimum $n_Q$ is nonexceptional over $X$.
For small rational $\delta>0$, the normalized divisor $D_s$ absorbs
$\delta f^*B$ at every nonexceptional prime.
Thus effectivity of $N+\delta B$ supplies a section on $X$.
Otherwise this pseudoeffective divisor is not big, so a nonzero movable
class $\alpha$ annihilates both $N$ and $B$.
The correction $\Theta$ of Proposition~\ref{g:bt:toroidal-transfer} is
supported on $B$.  Proposition~\ref{g:bt:unweighted-slope} gives
$\alpha\cdot M\leq0$ for every cotangent line $M$.

The determinant argument produces effective
$D_j\sim M+jN+B$, $D_j\geq B$, for unbounded $j$; their degree against
$\alpha$ is zero.  Resolve the Iitaka map of one $D_b$ after fixing
an earlier index $i<b$.  The same interpolation identity as above puts
every later restricted divisor in the support of $D_{b,S}$.
Include the finitely many exceptional divisors of the chosen resolution
in an auxiliary effective divisor $A$ and apply
Lemma~\ref{g:ba:rigidity} to $u^*D_b+A$.
Choose comparable indices $j_1<j_2$ in the tail beyond $b$, using
the ascending-pair property in a finite product of $\mathbb N$.
Put $v=j_2-j_1>0$. Then
\[
 C=(D_{j_2}-D_{j_1})|_S\geq0,\qquad C\sim vN_S,
 \qquad \operatorname{Supp}C\subseteq\operatorname{Supp}D_{b,S}.
\]
Let $\zeta$ be the corresponding section of $vN_S$, and let
$\sigma$ denote the normalized regular section of
$a\pi^*L-g^*N$ on a common resolution. The section
$\sigma^v g^*\zeta$ gives a new invariant generic section $s'$
of degree $avL$.

We verify the new ranks at exactly these multiples. Let $q\geq0$
and let $t$ be an invariant section of $K_{Z'/X}+qav\pi'^*L$
on the new generic fibre, where $Z'$ is a common smooth resolution.
The old rank condition, applied over $K$, says that division by
the Jacobian section and $\sigma^{qv}$ leaves the pullback of a
rational section $r$ of $qvN_S$. At a prime of $S$ outside $B_S$
and the $u$-exceptional locus, a nonexceptional minimizing component
has both normalized section order and Jacobian order zero.
Regularity of $t$ therefore makes $r$ regular at that prime.

The rational function $r/\zeta^q$ has poles only in
$\operatorname{Supp}(u^*D_b+A)|_S$: the divisor $D_b$ contains $B$,
the zeros of $\zeta$ lie in $D_{b,S}$, and $A$ contains every
$u$-exceptional prime of the fixed resolution.
Every multiple of this effective divisor is rigid by
Lemma~\ref{g:ba:rigidity}, so the function is constant over
$\mathbb C(Y_2)$. Thus $t$ is a base multiple of the Jacobian
times $(s')^q$. This also treats $q=0$, when $\zeta^q=1$.
Those Jacobian products supply nonzero sections in every tested
degree, proving rank exactly one. The argument permits extra
resolutions without assuming that their exceptional centres lie
in the original base locus.
\end{proof}

\subsection{Contraction using a large adjoint moment}
\label{g:ba:eventual}

\begin{proposition}\label{g:ba:eventual-contraction}
Suppose $K,s,d$ satisfy
\eqref{g:bt:horizontal-regular}--\eqref{g:bt:eventual-rank}.
If $\operatorname{trdeg}_{\mathbb C}K>0$, either a positive multiple of
$L$ has a nonzero invariant section, or these same two conditions hold
for some $K'\subset K$ with strictly smaller transcendence degree and a
new invariant rational section of a positive multiple of $L$.
\end{proposition}

\begin{proof}
Use Proposition~\ref{g:bt:large-moment} and put $N^+=N+R/k$.
If $N^+$ is rationally effective, its section returns to $X$.
Otherwise $N^+$ is pseudoeffective and not big.  Choose a nonzero
movable $\alpha$ annihilating it.  Since $N,R$ are pseudoeffective,
\[
                         \alpha\cdot N=\alpha\cdot R=0.
\]
The unweighted metric of $J=-K_Y+kN+R$ and
Proposition~\ref{g:bt:unweighted-slope} give the cotangent bound for
$\alpha$.  The determinant construction supplies effective
$D_j\sim M+m_jN$ with $m_j$ increasing and $\alpha\cdot D_j=0$.

Let $K'$ be the field of functions with poles only on $\alpha$-null
primes.  Lemma~\ref{g:cf:null-field} gives finite generation, relative
algebraic closure, and strict transcendence-degree decrease.
On a resolution of the resulting map there are only finitely many
possible horizontal components of the $D_j$: the null horizontal primes
are finite because $\operatorname{Pic}(Y)\otimes\mathbb Q$ is finite
dimensional, and the resolution has only finitely many exceptional
primes.  The finite-support argument in the proof of Proposition~\ref{g:cf:dickson}
therefore gives $j>i$, $v=m_j-m_i>0$ and $c\in K^*$ with
\[
             D_j-D_i=vN+\operatorname{div}(c)
\]
horizontally effective over $K'$.
Set $s'=s^vc$.  Its divisor upstairs is the sum of the pullback of this
difference and $vD\geq0$.  Hence it is regular at every valuation
trivial on $K'$, proving \eqref{g:bt:horizontal-regular}.

For all sufficiently large $m$, a section $t$ satisfying the new
discrepancy inequalities satisfies the old ones, so
$t=s^{mv}u$ with $u\in K$.  If $Q$ has positive $\alpha$-degree,
$r_Q=0$.  A component attaining the two minima in the definition of
$R$ consequently has $a_E=0$ and $h_E=e_En_Q$.
Its valuation is trivial on $K'$, and the inequality for $t$ gives
\[
                       \operatorname{ord}_Q u+mv n_Q\geq0.
\]
The null support of $D_j-D_i$ gives
$\operatorname{ord}_Q c+v n_Q=0$.
Thus $u/c^m$ has no pole on any positive-degree prime and belongs to
$K'$.  We obtain $t\in K'(s')^m$.
The opposite inclusion follows from the regularity of $s'$ and the
nonnegative discrepancies over smooth $X$.
\end{proof}

\subsection{Original horizontal primes and the signed boundary}
\label{g:ba:signed}

The next variant tests poles only on primes of $X$ itself.
It therefore permits arbitrary poles at exceptional valuations in the
rank test, while the distinguished section remains regular on the
complete generic fibre.

\begin{proposition}\label{g:ba:original-primes}
Let $K=\mathbb C(Y)\subseteq\mathbb C(X)^T$ be relatively algebraically
closed, and let $s$ be an invariant rational section of $aL$, $a>0$, regular on
a resolved generic fibre over $K$.  Suppose that, for every $j\geq0$,
\begin{equation}\label{g:ba:original-rank}
 u\in\mathbb C(X)^T,\quad
 s^ju\text{ has no poles on primes of $X$ horizontal over $K$}
                  \quad\Longrightarrow\quad u\in K.
\end{equation}
Then either a positive multiple of $L$ has a nonzero invariant section,
or a contraction to a smaller base preserves these two conditions.
\end{proposition}

\begin{proof}
On a prepared equidimensional toroidal model $p:W\to X$, $f:W\to Y$,
normalize $s$ by vertical minima:
\[
 D_s=\operatorname{div}(s)-f^*N\geq0.
\]
The hypothesis is the fixed-section uniqueness condition of
Lemma~\ref{g:cf:uniqueness}. It gives invariant adjoint rank one
in every nonnegative multiple of $a$: divide by the Jacobian and
test the nonexceptional horizontal primes. Hence
Lemma~\ref{g:cf:metric-conclusions} supplies a bounded semipositive
metric on $N$ and a signed boundary $B_{\mathrm{sgn}}$ with simple
normal crossing support and coefficients less than one, such that
$-(K_Y+B_{\mathrm{sgn}})$ is nef.

We compare the two choices of permissible base poles before choosing
the contraction. Let $S_0$ be the reduced sum of omitted primes,
above which every component is exceptional over $X$. Define $N'$
by taking the minimum only over nonexceptional components off $S_0$,
and retaining the coefficient of $N$ on $S_0$. Put $G=N'-N\geq0$.
Let $B_{\min}$ be the reduced sum of primes at which no minimizing
component is nonexceptional. Then
\[
 \operatorname{Supp}B_{\min}
   =\operatorname{Supp}S_0\ \sqcup\ \operatorname{Supp}G,
 \qquad
 \operatorname{Supp}B_{\mathrm{sgn}}^-
   \subseteq\operatorname{Supp}S_0.
\]
Indeed, away from the omitted primes, failure to attain the full
minimum on a nonexceptional component is exactly the strict inequality
$n'_Q>n_Q$.

There are two effectivity tests. If $N'+S_0$ is rationally effective,
its section returns to $X$ by the nonexceptional minimum. Alternatively,
for sufficiently small rational $\epsilon>0$,
$D_s-\epsilon f^*B_{\min}$ has no negative coefficient at any
nonexceptional prime; effectivity of $N+\epsilon B_{\min}$ therefore
also returns a section. If neither test succeeds, the chosen
pseudoeffective class is not big. Cone duality gives a nonzero movable
class $\alpha$. Either choice yields the same vanishing conditions:
\[
 \alpha\cdot N=\alpha\cdot G=\alpha\cdot S_0=0,
 \quad\text{equivalently}\quad
 \alpha\cdot N=\alpha\cdot B_{\min}=0.
\]
In particular $\alpha\cdot B_{\mathrm{sgn}}^-=0$, so
Proposition~\ref{g:bt:signed-slope} supplies the cotangent bound.
The determinant construction gives effective
$D_j\sim M+jN$ at unbounded divisible indices, with
$\alpha\cdot D_j=0$.

Let $K'$ be the null-prime field of Lemma~\ref{g:cf:null-field};
it has smaller transcendence degree. On a common resolution, the
affine interpolation of Lemma~\ref{g:cf:horizontal-slope} gives
$j>i$ and
\[
 D_j-D_i=vN+\operatorname{div}(c),\qquad v=j-i>0,\quad c\in K^*,
\]
with effective horizontal part over $K'$. One may instead use
the finite-support extraction of Proposition~\ref{g:cf:dickson}
to choose a comparable pair of horizontal multiplicity vectors.
Both choices give the same required property. Since the divisor
of $s'=s^vc$ is $vD_s+f^*(D_j-D_i)$, this section is regular on
the new complete generic fibre, including at exceptional horizontal
primes.

For uniqueness, suppose that $t=(s')^kg$ has no poles at original
primes horizontal over $K'$, where $k\geq0$ and
$g\in\mathbb C(X)^T$. Since these include the $K$-horizontal primes,
the old test applied to $t=s^{vk}c^kg$ gives $c^kg\in K$ and
hence $g\in K$. If a base prime $Q$ has positive $\alpha$-degree,
then $Q$ is not omitted, $n'_Q=n_Q$, and its valuation is trivial
on $K'$. Its nonexceptional components are therefore original
primes horizontal over $K'$. Testing $t$ on all those components
gives
\[
 \operatorname{ord}_Q g+
       k(vn'_Q+\operatorname{ord}_Q c)\geq0.
\]
The null support of $D_j-D_i$ gives
$vn_Q+\operatorname{ord}_Q c=0$. Thus $g$ has no pole at any
positive-degree prime. Its pole divisor is null, and the field
criterion in Lemma~\ref{g:cf:null-field} gives $g\in K'$.
This includes $k=0$ and proves the required uniqueness.

This proof uses the whole bad-minimum set through $G$ and $S_0$.
Annihilating only $B_{\mathrm{sgn}}^-$ would not provide the
nonexceptional minimizing components needed in the last pole test.
\end{proof}

\subsection{A toric construction of the initial generic section}
\label{g:ba:toric-start}

The rational torus quotient also admits a geometric construction of its
first invariant section, separate from wedging fundamental vector
fields.  This explains why a generated line on a generic orbit
compactification suffices, without assuming semiampleness on $X$.

\begin{proposition}\label{g:ba:toric-section}
For the rational quotient $K=\mathbb C(X)^T$, a positive multiple of
$L$ has a nonzero invariant section on the smooth resolved generic
fibre.  The data satisfy \eqref{g:ba:original-rank}.
\end{proposition}

\begin{proof}
Smooth semipositivity implies nefness, so
Lemma~\ref{g:cf:toric-initial-section} supplies the invariant section.
Ratios of invariant rational tensors of the same power of \(L\) lie in
\(\mathbb C(X)^T=K\), which proves \eqref{g:ba:original-rank}.
\end{proof}

The preceding contractions preserve their stated data and strictly
decrease a nonnegative base dimension.  At a point base the generic
section is global and pushes to $X$.  They therefore give complete
alternative reductions to an invariant section under their respective
initial conditions.
